\label{sec:tests}

The \textsc{bisect-apportion} test suite comprises 61 tests organised
into three levels following the \textsc{bisect} platform convention.

\subsection{L0: Unit Tests (37 tests)}

L0 tests verify individual functions and edge cases in isolation.

\paragraph{Priority formula tests (12 tests):}
\begin{itemize}
  \item \texttt{hh\_priority(100, 1)}: Expected \texttt{Priority\{num: 10000, den: 2\}}
    (i.e., $100^2/(1\times2) = 5000$, representing priority $\approx 70.7$).
  \item \texttt{hh\_priority(100, 2)}: Expected \texttt{Priority\{num: 10000, den: 6\}}
    ($100^2/(2\times3) \approx 1666.7$).
  \item \texttt{webster\_priority(100, 0)}: $P = 2\times100/1 = 200$.
  \item \texttt{adams\_priority(100, 1)}: $P = 100/1 = 100$.
  \item \texttt{jefferson\_priority(100, 0)}: $P = 100/(0+1) = 100$.
  \item And corresponding tests for Webster, Adams, Jefferson at seats 1--3.
\end{itemize}

\paragraph{Comparison tests (8 tests):}
\begin{itemize}
  \item Two equal-priority states compared exactly (cross-multiplication).
  \item Priority ordering: HH threshold $< $ Webster threshold for $s=1$.
  \item Priority comparison: states with populations 100 and 101 at seat 1.
\end{itemize}

\paragraph{Edge case tests (17 tests):}
\begin{itemize}
  \item Single state, $H=1$: returns [(``State'', 1)].
  \item Single state, $H=10$: returns [(``State'', 10)].
  \item Two states, equal population, $H=4$: returns [(``A'', 2), (``B'', 2)].
  \item $H = n$ (minimum House): all states receive exactly 1 seat.
  \item State with population 1 and state with population $10^6$: large
    population state gets vastly more seats.
  \item Constitutional floor: state with very small population
    receives 1 seat (not 0).
  \item Hamilton at exact quotas: no fractional parts, same as lower quotas.
  \item \texttt{check\_alabama\_paradox}: Hamilton on 3-state paradox
    configuration returns \texttt{Some(...)}.
  \item \texttt{check\_alabama\_paradox}: HH on 2020 Census returns \texttt{None}.
\end{itemize}

\subsection{L1: Integration Tests (14 tests)}

L1 tests verify complete apportionments on synthetic examples with
known correct answers.

\paragraph{5-state canonical example (5 tests):}
Populations $(6, 6, 2, 2, 1)$, $H = 10$:
\begin{itemize}
  \item HH: verify hand-computed priority sequence produces expected allocation.
  \item Webster: expected same result for this example.
  \item Adams: expected same result.
  \item Jefferson: expected result (floor rounding, possibly different).
  \item Hamilton: expected result (exact quotas happen to be near integers).
\end{itemize}

\paragraph{Paradox detection (3 tests):}
\begin{itemize}
  \item 3-state Alabama paradox configuration:
    \texttt{check\_alabama\_paradox} returns the correct state name.
  \item HH on same configuration: \texttt{None} (paradox-free).
  \item Hamilton on additional paradox configuration (from J.5): correct state.
\end{itemize}

\paragraph{Verification protocol tests (3 tests):}
\begin{itemize}
  \item \texttt{verify\_sha256} returns \texttt{true} for correct allocation.
  \item \texttt{verify\_sha256} returns \texttt{false} for off-by-one allocation.
  \item Canonical string format matches expected output for 3-state example.
\end{itemize}

\paragraph{D'Hondt equivalence test (1 test):}
5-party example from \citet{lijphart1994electoral}: Jefferson with vote totals
produces same allocation as standard d'Hondt table.

\paragraph{Webster-Sainte-Lagu\"{e} equivalence test (1 test):}
5-party example with Webster: produces same allocation as Sainte-Lagu\"{e}.

\paragraph{Overflow stress test (1 test):}
Populations at maximum realistic values ($5 \times 10^7$), 52 seats:
no \texttt{u128} overflow occurs.

\subsection{L2: Real-Data Regression Tests (10 tests)}

L2 tests verify against Census Bureau official apportionments.

\paragraph{2020 Census (3 tests):}
\begin{itemize}
  \item HH 2020: SHA verification passes; all 50 states match official table.
  \item Webster 2020: all 50 states match expected Webster output.
  \item Hamilton 2020: all 50 states match expected Hamilton output
    (which agrees with HH for 2020 data).
\end{itemize}

\paragraph{2010 Census (3 tests):}
\begin{itemize}
  \item HH 2010: SHA verification passes; all 50 states match.
  \item Webster 2010: matches expected (may differ from HH for 1--2 states).
  \item Hamilton 2010: matches expected.
\end{itemize}

\paragraph{2000 Census (3 tests):}
\begin{itemize}
  \item HH 2000: SHA verification passes; all 50 states match.
  \item Webster 2000: matches expected.
  \item Adams 2000: matches expected.
\end{itemize}

\paragraph{Historical census test (1 test):}
HH for 1940 Census (the last apportionment under Webster before 1941):
produces the officially published 1940 apportionment.
